\chapter{模版的使用}\label{ch:mannual}
本手册使用\LaTeX{}编译生成，本章的例子可以找到对应的TEX源代码。

\section{封面的生成}

非常重要！！！

由于本人能力和时间有限，所以没能用\LaTeX 做成学校论文的中英文封面和原创性声明，所以需要使用的用户通过学校提供的模板手工生成pdf文件，然后再通过导入的方法使用。
而且需要注意的是，考虑到打印的问题，请在中文封面，英文封面和原创性声明之后都手工加入一个空白页(Ctrl+Enter)。
\section{dmuthesis.tex主源文件}

\subsection{论文封面及原创性声明}
dmuthesis.tex是本模版的主源文件，所有的章节都需通过这个文件包含进来，也是Xe\LaTeX{}命令编译的对象。
这个文件的主要语句使用如下：

修改\verb|/cover/thesis_cover.doc|，转化成pdf文件，用以下命令包含论文封面：
\begin{verbatim}
\includepdf[pages=-]{cover/thesis_cover.pdf}
\end{verbatim}
使用以上方法，硕士也可以使用本模板进行写作了。
\subsection{修改论文题目}
\begin{verbatim}
\newcommand{\thesissubject}{大连海事大学研究生论文模板使用帮助}
\end{verbatim}

\subsection{摘要与目录}
由于硕士论文不要求创新点摘要，所以硕士研究生请将下面两行代码注释。
\begin{verbatim}
\originalabstract 
\begin{abstractOriginal}
豫章故郡，洪都新府。星分翼轸，地接衡庐。襟三江而带五湖，控蛮荆而引瓯越。物华天宝，龙光射牛斗之墟；人杰地灵，徐孺下陈蕃之榻。雄州雾列，俊彩星驰。台隍枕夷夏之交，宾主尽东南之美。都督阎公之雅望，綮戟遥临；宇文新州之懿范，譫帷暂驻。十旬休暇，胜友如云。千里逢迎，高朋满座。腾蛟起凤，孟学士之词宗；紫电青霜，王将军之武库。家君作宰，路出名区。童子何知？躬逢胜饯。

时维九月，序属三秋。潦水尽而寒潭清，烟光凝而暮山紫。俨骖騑于上路，访风景于崇阿。临帝子之长洲，得仙人之旧馆。层峦耸翠，上出重霄；飞阁流丹，下临无地。鹤汀凫渚，穷岛屿之萦回；桂殿兰宫，即冈峦之体势。

披绣闼，俯雕甍。山原旷其盈视，川泽纡其骇瞩。闾阎扑地，钟鸣鼎食之家；舸舰迷津，青雀黄龙之舳。虹销雨霁，彩彻区明。落霞与孤鹜齐飞，秋水共长天一色。渔舟唱晚，响穷彭蠡之滨，雁阵惊寒，声断衡阳之浦。

遥吟甫畅，逸兴遄飞。爽籁发而清风生，纤歌凝而白云遏。睢园绿竹，气凌彭泽之樽；邺水朱华，光照临川之笔。四美具，二难并。穷睇眄于中天，极娱游于暇日。天高地迥，觉宇宙之无穷；兴尽悲来，识盈虚之有数。望长安于日下，指吴会于云间。地势极而南溟深，天柱高而北辰远。关山难越，谁悲失路之人；萍水相逢，尽是他乡之客。怀帝阍而不见，奉宣室以何年？

嗟乎！时运不齐，命途多舛。冯唐易老，李广难封。屈贾谊于长沙，非无圣主；窜梁鸿于海曲，岂乏明时？所赖君子安贫，达人知命。老当益壮，宁移白首之心；穷且益坚，不坠青云之志。酌贪泉而觉爽，处涸辙而犹欢。北海虽赊，夫摇可接；东隅已逝，桑榆非晚。孟尝高洁，空怀报国之情；阮籍猖狂，岂效穷途之哭？

\end{abstractOriginal}



\end{verbatim}

由于学校模板中并没有要求加入图目录和表格目录，所以这里也没有加入。

\section{新建章节}

正文的章节使用\verb|\chapter|，例如：\verb|\chapter{绪论}|

新建一章的第一级小节命令：\verb|\section{新建章节}|。
多级章小节用\verb|\subsection{3.2.1}|，\verb|\subsubsection{3.2.1.1}|。不建议使用超过4级的小节。若有需要可以使用没有编号的章节题目。使用`*'去掉编号，命令\verb|\subsubsection*{无编号章节}|，可以得到以下效果：

\subsubsection*{无编号章节}

参考文献之后的章节用\verb|\appendix|命令，如\verb|\appendix{致谢}|。

\section{字体控制}

模板预制了多种字体，可供参考，比如{\songti 宋体}，{\heiti 黑体 },{\kaiti 楷体}和{\fangsong 仿宋}。

字号也有很多可供选择，比如 

{\chuhao 初号} 就是太大了。

{\xiaochuhao 小初号} 也非常大 。

{\yichu 一初} 我也不知道为什么要设置这个字号。
 
{\yihao 一号} 用的地方也比较少。
  
{\erhao 二号}。

{\xiaoerhao 小二号}。

{\sanhao 三号}。

{\xiaosanhao 小三号}。

{\sihao 四号} 和{\xiaosihao 小四号} 。

{\wuhao 五号}和 {\xiaowuhao 小五号}。
 
{\liuhao 为什么你要用六号字呢？} 

{\xiaoliuhao 用小六号，你是心虚不想让别人看清楚吧} 。

{\qihao 七号，你在搞什么啊？} 

除此之外，还可以加入不同的{\color {blue} 颜}  {\color {cyan} 色}。{\xiaoerhao {\color{blue}配合} \sihao 不同的\kaiti 字体\fangsong  和字号，{\color{yellow}如果}你真的\heiti {\color{red}有时间的话}}。


\textbf{所以，这些字号的设置啊，就让它随波而去吧，提供了不代表你要使用！}
\section{插入图片}
\subsection{一张图片}

\begin{figure}[H]
\centering
\includegraphics[scale=0.4]{figs/dmu-logo.eps}
\FigureBicaption{大连海事大学标志}{DMU LOGO}
\label{fig:single}
\end{figure}


以下命令将产生图\ref{fig:single}的效果：
\begin{verbatim}
\begin{figure}[H]
\centering
\includegraphics[scale=0.4]{figs/dmu-logo.eps}
\FigureBicaption{大连海事大学标志}{DMU LOGO}
\label{fig:single}
\end{figure}
\end{verbatim}
其中[H]参数强制固定浮动图形的位置；scale参数可以调整图片大小；\verb|\FigureBicaption{中文标题}{英文标题}|加入图片标题；\verb|\label|命令用来引用图片，具体用法见第\ref{sect:ref}节。
建议使用\verb|width=0.7\textwidth|代替\verb|scale=0.4|，可以更方便控制图片大小。

\subsection{并列子图}\label{subsect:subfigeres}


\begin{figure}
\begin{minipage}{0.48\textwidth}
  \centering
  \includegraphics[height=1.5cm]{figs/dmu-logo.eps}
  \FigureBicaption{并排第一个图}{Two parallel figures}
  \label{fig:parallel1}
\end{minipage}\hfill
\begin{minipage}{0.48\textwidth}
  \centering
  \includegraphics[height=1.5cm]{figs/dmu-logo.eps}
  \FigureBicaption{并排第二个图}{Two parallel figures}
  \label{fig:parallel2}
\end{minipage}
\end{figure}
如果多个图形相互独立，并不共用一个图形计数器，用 \verb|minipage|。以下命令效果见图\ref{fig:parallel1}和图\ref{fig:parallel2}：
\begin{verbatim}
\begin{figure}
\begin{minipage}{0.48\textwidth}
  \centering
  \includegraphics[height=2cm]{figs/dmu-logo.eps}
  \FigureBicaption{并排第一个图}{Two parallel figures}
  \label{fig:parallel1}
\end{minipage}\hfill
\begin{minipage}{0.48\textwidth}
  \centering
  \includegraphics[height=2cm]{figs/dmu-logo.eps}
  \FigureBicaption{并排第二个图}{Two parallel figures}
  \label{fig:parallel2}
\end{minipage}
\end{figure}
\end{verbatim}

多个子图的情况用见图\ref{fig:subfigs-singleCaption}，这个图中，子图只有中文的标题。
\begin{figure} %[h]
  \centering%
  \subfloat[第一个小图形]{%
    \label{fig:subfig1-singleCaption}
    \includegraphics[height=2cm]{figs/dmu-logo.eps}}\hspace{4em}%
  \subfloat[第二个小图形。如果标题很长的话，它会自动换行，这个 caption 就是这样的例子]{%
    \label{fig:subfig2-singleCaption}
    \includegraphics[height=2cm]{figs/dmu-logo.eps}}
  \FigureBicaption{多个子图}{Subfigures}
  \label{fig:subfigs-singleCaption}
\end{figure}

子图也使用双标题，其效果见图\ref{fig:subfigs}
\begin{figure} %[h]
  \centering%
  \SubfloatCaption{子图中文标题}{English caption}{fig:subfig1}{\includegraphics[height=2cm]{figs/dmu-logo.eps}}
  \hspace{4em}%
  \SubfloatCaption{子图中文标题}{English caption}{fig:subfig2}{\includegraphics[height=2cm]{figs/dmu-logo.eps}}
  \FigureBicaption{多个子图}{Subfigures}
  \label{fig:subfigs}
\end{figure}
\begin{verbatim}
\begin{figure} %[h]
  \centering%
  \SubfloatCaption{子图中文标题}{English caption}{fig:subfig1}
    {\includegraphics[height=2cm]{figs/dmu-logo.eps}}
  \hspace{4em}%
  \SubfloatCaption{子图中文标题}{English caption}{fig:subfig2}
    {\includegraphics[height=2cm]{figs/dmu-logo.eps}}
  \FigureBicaption{多个子图}{Subfigures}
  \label{fig:subfigs}
\end{figure}
\end{verbatim}

\section{表格}
\subsection{基本表格}
\label{subsect:basictable}
双标题表格，由于学校模板中并没有对子表格规定是否要用双标题，所以请自行决定。

\begin{table}[htb]
  \centering
  \begin{minipage}[t]{0.8\linewidth} % 如果想在表格中使用脚注，minipage是个不错的办法
  \TableBicaption{如果表格的标题很长，那么在表格索引中就会很不美
    观，所以要像 chapter 那样在前面用中括号写一个简短的标题。这个标题会出现在索
    引中。}{English caption}
  \label{tab:basic}
    \begin{tabular*}{\linewidth}{lp{10cm}}
      \toprule[1.5pt]
      {文件名} & {描述} \\\midrule[1pt]
      scutthesis.cls & 模板类文件。\\
      scutthesis.bst    & 参考文献 Bibtex 样式文件。\\
      \bottomrule[1.5pt]
    \end{tabular*}
  \end{minipage}
\end{table}


\subsection{并列表格/子表格}
与并列图形方法类似,参考第\ref{subsect:subfigeres}节。
\begin{table}%[h]
\noindent\begin{minipage}{0.5\textwidth}
\centering
\caption{第一个并排子表格}
\label{tab:parallel1}
\begin{tabular}{p{2cm}p{2cm}}
\toprule[1.5pt]
111 & 222 \\\midrule[1pt]
222 & 333 \\\bottomrule[1.5pt]
\end{tabular}
\end{minipage}
\begin{minipage}{0.5\textwidth}
\centering
\caption{第二个并排子表格}
\label{tab:parallel2}
\begin{tabular}{p{2cm}p{2cm}}
\toprule[1.5pt]
111 & 222 \\\midrule[1pt]
222 & 333 \\\bottomrule[1.5pt]
\end{tabular}
\end{minipage}
\end{table}

\begin{table}
\centering
\caption{并排子表格}
\label{tab:subtable}
\subfloat[第一个子表格]{
\begin{tabular}{p{2cm}p{2cm}}
\toprule[1.5pt]
111 & 222 \\\midrule[1pt]
222 & 333 \\\bottomrule[1.5pt]
\end{tabular}}\hskip2cm
\subfloat[第二个子表格]{
\begin{tabular}{p{2cm}p{2cm}}
\toprule[1.5pt]
111 & 222 \\\midrule[1pt]
222 & 333 \\\bottomrule[1.5pt]
\end{tabular}}
\end{table}

模版还提供了更加复杂的表格功能，如表格中的斜线，注释等。本文档暂时不提供这些复杂表格的例子，暂留给读者探索。跨页长表格的例子可以看到附录中的发表文文章列表。

\section{公式}
简单公式环境：
\begin{equation}\label{eq:simple}
y=mx+c
\end{equation}

公式太长，多行排列：
\begin{equation}\label{eq:split}
  \begin{split}
      y&=mx+c\\
       & \quad(n+o)x+c\\
       &=1
  \end{split}
\end{equation}
另外一种方法：
\begin{multline}
1+2+3+4+\\
5+6+7+8
\end{multline}
多个公式并列不要用多个equation环境，会造成公式间距过大的问题，用gather环境：
\begin{gather}
y=mx+c \label{eq:eq1}\\
x=(n-2)+d \label{eq:eq2}
\end{gather}
公式的对齐，以等号对齐为例：
\begin{align}
a&=b+c\\
d&=e+f+g+h
\end{align}
再来一个例子：
\begin{align}
a{11}&=1 & b_{12}&=2\\
a{21}&=3 & b_{22}&=4
\end{align}
来一个对齐方式不一样的：
\begin{flalign}
a{11}&=1 & b_{12}&=2\\
a{21}&=3 & b_{22}&=4
\end{flalign}

\section{定理环境}
模版提供了多种定理环境：命题(proposition)，引理(lemma)，定理(theorem)，公理(axiom)， 推论(corollary)，情形(case)，猜想(conjecture)，性质(property)，还有定义(definition)，例(example)，注(remark)。下面以常见的环境作例子。\cite{Barman2006,Berry2007,Boutsidis2008,Buciu2006,Cai2005,Chen2008,Cichocki2006,Cichocki2008,Cichocki2007,Cichocki2007a,Cichocki}
\begin{theorem}[人论]\label{thm:ex1}
你是人。
\end{theorem}
把证明内容放到\verb|\begin{proof}|和\verb|\end{proof}|即可得到：
\begin{proof}
是的，大家都说你有点像人$\sim\sim$
\end{proof}

\section{列表环境}
无编号列表
\begin{itemize}
\item 第一项
	\begin{enumerate}
	\item 下一次
	\item 在下一次
	\end{enumerate}
\item 第二项
\end{itemize}

编号列表
\begin{enumerate}
\item 第一项
\item 第二项
\end{enumerate}

描述环境，一般用于几个定义的描述：
\begin{description}
  \item[天才] 不按常理思考；
  \item[流氓] 不按道德行动；
\end{description}

\section{参考文献}

关于参考文献的管理以及使用，请参考 \url{http://aurora1625.github.com/blog/2012/02/16/tame-the-beast/}。此处不再赘述。


\section{算法}\label{sect:algNcode}
伪代码可以用algorithm环境，下面的代码得到算法\ref{alg: life}的效果。
\begin{verbatim}
\begin{algorithm}
\caption{生命的轨迹} \label{alg: life}
\nl 睁眼\;
\nl \If {O了没？} {O了，闭眼}
\nl \Else{闭眼吧，忙活了一辈子，这时候还能有什么放不下。}
\end{algorithm}
\end{verbatim}
\begin{algorithm}
\caption{生命的轨迹} \label{alg: life}
\nl 睁眼\;
\nl \If {O了没？} {O了，闭眼}
\nl \Else{闭眼吧，忙活了一辈子，这时候还能有什么放不下。}
\end{algorithm}
%
%源代码用lstlisting环境。把代码\ref{code:love}的内容放在\verb+\begin{lstlisting}+和
%\verb+\end{lstlisting}+即可得到代码\ref{code:love}的效果
%\begin{lstlisting}[language={C}, label=code:love,caption={恋爱程序}]
%include<love.h>
%char result love(boy, girl)
%　　　　{
%　　　　    if( boy.有房() and boy.有车() )
%　　　　    {
%　　　　    boy.set(nothing);
%　　　      return girl.嫁给(boy);
%　　　　    }
%　　　　    if( girl.愿意等() )
%　　　　    {
%　　　　    while(!(boy.赚钱 > 100,000 and girl.感情 > 8 ) )
%　　　　    {
%　　　　    for( day=1; day <=365; day++)
%　　　　    {
%　　　　        if( day == 情人节 )
%　　　　        if( boy.givegirl(玫瑰) )
%　　　　        girl.感情++;
%　　　　        else
%　　　　        girl.感情--; 　　
%　　　　        if( day == girl.生日)
%　　　　        if( boy.givegirl(玫瑰) )
%　　　　        girl.感情++;
%　　　　        else
%　　　　        girl.感情--;
%　　　　        boy.拼命赚钱();
%　　　　     }
%　　　　 }
%　　　　　　　 if( boy.有房() and boy.有车() )
%　　　　　　　 {
%　　　　          boy.set(nothing);
%　　　　          return girl.嫁给(boy);
%　　　　      }
%　　　　      年龄++;
%　　　　      girl.感情--;
%　　　　      }
%　　　　      return girl.goto( another_boy);
%　　　　 }
%\end{lstlisting}
\section{交叉引用}\label{sect:ref}
交叉引用需要两个步骤。
\begin{enumerate}
 \item 用\verb|\label{}|命令标识；
 \item 用\verb|\ref{}|命令引用。
\end{enumerate}

从本节的例子可以看出，无论是图片，表格，公式，定理，算法，代码，章节等，都可以表示和引用。如\verb|\label{ch:intr}|和\verb|第\ref{ch:intr}章|。注意到，\verb|\ref|命令只是引用了编号，并没有给出引用的类型，因此需要加上引用类型的名字，再如\verb|算法\ref{alg: life}|。公式的编号一般在括号里，特殊地，可以用\verb|\eqref{}|代替`\verb|(\ref{})|'。



